\section{Higher Morita categories and invertibility}
\label{sec:Morita}

\subsection{Conventions}
\label{sec:conventions}

Throughout, $k$ is a field of characteristic zero. We write $\Pre$ for the symmetric monoidal
$\left(\infty,2\right)$-category of presentable stable $k$-linear $\infty$-categories,
colimit-preserving $k$-linear functors and natural transformations, with the Lurie tensor product.
A stack is a derived stack over $k$; further hypotheses (perfect, QCA, bounded) are stated where
they are used. Formal completions and $\IndCoh_0$ are recalled in Section \ref{sec:IndCoh}.
We write $d$ for the dimension in $\umor_d$ and $n$ for the shift in $\DMod^n$.

Furthermore, $V$ always denotes a vector space and $V^\vee$ its dual. Shifts are written
$V\left[n\right]$ with the convention that $V\left[1\right]\simeq BV$; in particular it is an odd vector space, and
$T^*\left[2\right]\left(V\left[1\right]\right)\simeq\left(V\dsum V^\vee\right)\left[1\right]$.
Duals are always written
$\left(-\right)^\vee$; we reserve the notation $\left(-\right)^*$ for pullbacks, cotangent bundles and
(co)chains.

\begin{rmk}[Abelian versus derived]
\label{rmk:abelian}
Much of the literature on braided tensor categories that we use, e.g.\
\cite{BJS, BJSS, Neg:frob}, works in the $\left(2,2\right)$-category of locally
presentable $k$-linear (abelian) categories, whereas we work with presentable stable
$k$-linear $\infty$-categories (Section \ref{sec:conventions}). We pass from the former to the
latter by taking derived categories, and we only do this for fusion categories, where nothing is
lost: A finite semisimple category with simple objects $I$ becomes the product 
$\prod_{i\in I} \Vect$, and relative tensor products over fusion categories have no higher
derived terms, since module categories over a fusion category are semisimple. For
non-semisimple categories the derived relative tensor products and centers can differ from
the abelian ones, and we cite results about them only as statements in the abelian
setting. Results imported from the abelian setting are marked by a reference to this
remark.
\end{rmk}

\subsection{Morita category of \texorpdfstring{$\EE_d$}{Ed}-algebras}
\label{sec:mor}

Let $\coef$ be a symmetric monoidal $\infty$-cate\-gory which admits sifted colimits,
preserved by $\otimes$ in each variable, as in \cite{SSS}.
The Morita $\left(\infty,d+1\right)$-category $\umor_d\left(\coef\right)$ of \cite{SSS} has
$\EE_d$-algebras as objects, bimodules in $\EE_{d-1}$-algebras as $1$-morphisms, bimodules of
bimodules in $\EE_{d-2}$-algebras as $2$-morphisms, and so on, with homomorphisms as
$\left(d+1\right)$-morphisms. It is symmetric monoidal, and its loops are the Morita category
of $\EE_{d-1}$-algebras,
\begin{equation*}
\Omega\,\umor_d\left(\coef\right)\simeq\umor_{d-1}\left(\coef\right) \ ,
\end{equation*}
cf.\ \cite[Proposition 3.3.8]{Sch:thesis}, where this is shown for the
factorization-algebra model of $\umor_d$.
For $\coef$ a (symmetric monoidal) $\left(\infty,m\right)$-category, one
obtains a (symmetric monoidal) $\left(\infty,d+m\right)$-category \cite{JFS}. In particular, since $\Pre$ is a symmetric monoidal
$\left(\infty,2\right)$-category, $\umor_d\left(\Pre\right)$ is a symmetric monoidal
$\left(\infty,d+2\right)$-category.

In our examples, such as $\coef=\Pre$ and $\coef=\Mod_R\left(\Pre\right)$ for
$R\in\CAlg\left(\Pre\right)$, $\coef$ is not presentable, but it admits all small colimits,
preserved by $\otimes$ in each variable \cite[Section 4.8.1]{L:HA}.

\subsection{Dualizability of \texorpdfstring{$\mathbb{E}_d$}{Ed}-algebras}
\label{sec:dual_Ek}

By taking opposite algebra structures in each of the $d$ directions, every $\EE_d$-algebra is $d$-dualizable in $\umor_d\left(\coef\right)$, see \cite{GS}. For TQFT applications, the more interesting question is
 when it is $\left(d+1\right)$-dualizable, i.e.\ when it defines a fully extended
framed $\left(d+1\right)$-dimensional TQFT.
This is answered by the following theorem, conjectured by Lurie
\cite[Remark 4.1.27]{L:Cob}.

\begin{theorem}[{\cite[Theorem A]{SSS}}]
\label{thm:SSS_dual}
An $\EE_d$-algebra $A$ in $\coef$ is $\left(d+1\right)$-dualizable in $\umor_d\left(\coef\right)$ if,
and only if, for all $0\leq i\leq d$, the object $A$ is left dualizable as a left module over
the factorization homology
\begin{equation*}
\int_{S^{i-1}\times\RR^{d-i+1}}A \ .
\end{equation*}
\end{theorem}

For $i=0$ the factorization homology is $\int_{\emptyset}A = \one_{\coef}$, so the condition is
dualizability of $A$ in $\coef$.
Theorem \ref{thm:SSS_dual} also applies to Morita categories valued in a symmetric monoidal
$\left(\infty,m\right)$-category such as $\Pre$, since $\left(d+1\right)$-dualizability is
detected in the underlying $\left(\infty,d+1\right)$-category \cite[Section 1.5]{SSS}.

\begin{rmk}
\label{rmk:SSS_relative}
There is also a relative version for $1$-morphisms $M\colon A\to B$ in
$\umor_d\left(\coef\right)$, characterizing when $M$ is $d$-times right adjointable in terms of
dualizability over factorization homologies of $M$ \cite[Theorem B]{SSS}. Via \cite{JFS}, such
$M$ defines an oplax relative field theory. Two cases are relevant below. The regular module
$A$ always gives a maximally adjointable $1$-morphism. An $A$-central $\EE_{d-1}$-algebra $N$,
i.e.\ a $1$-morphism $A\to\one$
(Notation \ref{nota:E2}(4)), is
maximally adjointable if, and only if, $N$ is $d$-dualizable as an object of
$\umor_{d-1}\left(\coef\right)$ \cite[Section 4.5]{SSS}. For $d=2$, this is the dualizability input for the boundary
conditions of Section \ref{sec:line_ops}.
\end{rmk}

\subsection{Invertibility of \texorpdfstring{$\mathbb{E}_d$}{Ed}-algebras}
\label{sec:inv_Ek}

For an $\EE_d$-algebra $A$ and $0\leq j\leq d$, write
$\mathcal{Z}_j\left(A\right)\ceqq\End_{\int_{S^{j-1}\times\RR^{d-j+1}}A}\left(A\right)$ for its
$\EE_j$-center (cf.\ Notation \ref{nota:E2}). The action of
$\int_{S^{i-1}\times\RR^{d-i+1}}A$ on $A$ induces a canonical map to $\mathcal{Z}_{d-i}\left(A\right)$.

\begin{cor}[{\cite[Corollary A]{SSS}}]
\label{cor:SSS_inv}
An $\EE_d$-algebra $A$ in $\coef$ is invertible in $\umor_d\left(\coef\right)$ if, and only if, it
is $\left(d+1\right)$-dualizable and, for all $0\leq i\leq d$, the canonical map
\begin{equation*}
\int_{S^{i-1}\times\RR^{d-i+1}}A\longrightarrow \mathcal{Z}_{d-i}\left(A\right)
\end{equation*}
to the $\EE_{d-i}$-center is an equivalence.
\end{cor}

In particular, \cite{SSS} proves \cite[Conjecture 1.9]{BJSS}.
For $d=2$, the conditions are those of Definition \ref{defn:inv_E2} below:
non-degeneracy ($i=0$), factorizability ($i=1$) and cofactorizability ($i=2$).

\subsubsection{Invertibility of \texorpdfstring{$\mathbb{E}_2$}{E2}-algebras}
\label{sec:inv_E2}

This section unpacks Corollary \ref{cor:SSS_inv} for $d=2$.
We will make use of the usual description of $\mathbb{E}_2$-algebras as locally constant
factorization algebras on $\left(0,1\right)^2 \simeq \mathbb{R}^2$.

\begin{nota}
\label{nota:E2}
Let $A$ be an $\mathbb{E}_2$-algebra in $\coef$.
\begin{enumerate}
\item We write $A^{\rev}$ and $A^{\sigma\rev}$ for the pushforwards of $A$ along the
reflections of $\left(0,1\right)^2$ in the first and second coordinate. Rotating by $\pi$
provides equivalences $A^{\rev}\simeq A^{\sigma\rev}$ of $\mathbb{E}_2$-algebras, and
$A^{\rev}$ is the dual of $A$ in $\umor_2\left(\coef\right)$ \cite{GS}.
\item We write
\begin{equation*}
A^e\ceqq \int_{S^0\times\mathbb{R}\times\mathbb{R}}A\simeq A\otimes A^{\rev}\qquad\text{and}\qquad
\mathrm{HC}\left(A\right):=\int_{S^1\times\mathbb{R}}A \ ,
\end{equation*}
where $S^0\times\mathbb{R}\subset\mathbb{R}$ and $S^1\times\mathbb{R}\subset\mathbb{R}^2$
carry the framings bounding a disk, and $A\otimes A^{\sigma\rev}$ is obtained in the same
way from $\mathbb{R}\times S^0\times\mathbb{R}$. These are $\mathbb{E}_1$-algebras acting on $A$ via the radial coordinate. The notation $\mathrm{HC}\left(A\right)$ follows \cite{BJSS}, after the Harish-Chandra category of a braided tensor category.
\item The $\mathbb{E}_j$-centers of Section \ref{sec:inv_Ek} of an $\mathbb{E}_2$-algebra
$A$ are explicitly:
\begin{equation*}
\mathcal{Z}_0\left(A\right) \ceqq \End\left(A\right)  \ ,
\qquad
\mathcal{Z}_1\left(A\right) \ceqq \End_{A^e}\left(A\right)  \ , \quad\text{and}
\qquad
\mathcal{Z}_2\left(A\right) \ceqq \End_{\HC\left(A\right)}\left(A\right) \ .
\end{equation*}
\item A 1-morphism $A\rightarrow B$ in $\umor_2\left(\coef\right)$ is an $\mathbb{E}_1$-algebra in
$\left(A , B\right)$-bimodules, equivalently an $\mathbb{E}_1$-algebra $C$ in $\coef$
together with a map $A\otimes B^{\sigma\rev}\rightarrow \mathcal{Z}_1\left(C\right)$.
We call such data an $\left(A,B\right)$-central algebra. In particular, a $1$-morphism $B\to\one$ is a
$B$-central algebra, i.e.\ an $\EE_1$-algebra $C$ with a central structure $B\to \mathcal{Z}_1\left(C\right)$.
So, as in \cite{BJSS, SSS}, a boundary condition, or relative theory, for $B$ is a
$1$-morphism $B\to\one$. The reformulation of bimodules in terms of central algebras is
\cite[Section 5.3.1]{L:HA}.

\item For an $\mathbb{E}_1$-algebra $C$ we write $C^{!}\ceqq \mathrm{Hom}_{C^e}\left(C,C^e\right)$
and $C^\vee \ceqq \Hom\left(C , \one_{\coef}\right)$ for the duals of $C$
as a $C^e$-module and as an object of $\coef$, whenever these exist.
\end{enumerate}
\end{nota}

\begin{theorem}[{\cite[Theorem 2.26]{BJSS}}]
\label{thm:trivial_Alg2}
Let $\coef$ be a closed symmetric monoidal $\left(\infty,2\right)$-category admitting geometric
realizations, e.g.\ $\coef=\Pre$.
Let $B$ be an $\EE_2$-algebra in $\coef$ and let $C$ be a $B$-central algebra, viewed as a
1-morphism $C\colon B\to\one_{\coef}$ in $\umor_2\left(\coef\right)$. Then $C$ is invertible if,
and only if, $C$ is 2-dualizable in $\umor_1\left(\coef\right)$, that is, dualizable as an object of
$\coef$ and as a $C^e$-module, and the following maps are equivalences:
\begin{enumerate}
    \item the evaluation map $C^{!}\otimes_B C\to C^e$;
    \item the map $B\to \mathcal{Z}_1\left(C\right)$ given by the $B$-central structure on $C$;
    \item the evaluation map $C^{\vee}\otimes_{C\otimes_B C^{\rev}}C\to \one_{\coef}$;
    \item the map $C^{\rev}\otimes_B C\to \mathrm{End}\left(C\right)$ given by the left and
    right actions of $C$ on itself.
\end{enumerate}
\end{theorem}

Invertibility of $A$ itself is the case of Theorem \ref{thm:trivial_Alg2} in which
$B = A\otimes A^{\sigma\rev}$ and $C = A$, viewed as the evaluation $1$-morphism
$A\otimes A^{\sigma\rev}\to\one_{\coef}$ of the duality between $A$ and $A^{\sigma\rev}$, and the four conditions then assemble into the following three.

\begin{defn}[\cite{BJSS}]
\label{defn:inv_E2}
An $\mathbb{E}_2$-algebra $A$ in $\coef$ is
\begin{enumerate}
\item \emph{non-degenerate} if the unit map $\one_{\coef}\to \mathcal{Z}_2\left(A\right)$ is an
equivalence;
\item \emph{factorizable} if the action map $A\otimes A^{\sigma\rev}\to \mathcal{Z}_1\left(A\right)$ is an equivalence;
\item \emph{cofactorizable} if the action map
$\HC\left(A\right)\to \End\left(A\right)=\mathcal{Z}_0\left(A\right)$ is an equivalence.
\end{enumerate}
\end{defn}

\begin{theorem}[{\cite[Theorems 1.2 and 1.4]{BJSS};
\cite[Theorem A and Corollary A]{SSS}}]
\label{thm:inv_E2}
Let $A$ be an $\mathbb{E}_2$-algebra in $\coef$. Then $A$ is $3$-dualizable in
$\umor_2\left(\coef\right)$ if, and only if, $A$ is dualizable as an object of $\coef$, as a left
$A^e$-module, and as a left $\mathrm{HC}\left(A\right)$-module. Moreover, $A$ is invertible in
$\umor_2\left(\coef\right)$ if, and only if, it is $3$-dualizable, non-degenerate, factorizable and
cofactorizable.
\end{theorem}

This is the case $d=2$ of Theorem \ref{thm:SSS_dual} and Corollary \ref{cor:SSS_inv}; cf.\
\cite{BJS} for related results on braided tensor categories in the abelian setting.

\subsubsection{Braided tensor categories}
\label{sec:braided}

Braided tensor categories are $\mathbb{E}_2$-algebras in the $\left(2,2\right)$-category of
locally presentable $k$-linear categories. The results below are stated there; we only use them
in $\Pre$ for fusion categories (Remark \ref{rmk:abelian}).
In either setting, invertibility in the
Morita category $\umor_2$ is detected in its underlying $\left(\infty , 3\right)$-category
\cite[Section 1.5]{SSS}.

In Definition \ref{defn:inv_E2}, $\mathcal{Z}_1\left(A\right)$ is the Drinfeld center and, in the abelian
setting, $\mathcal{Z}_2\left(A\right)$ is the M\"uger center of $A$ \cite[Propositions 3.6 and 3.7]{BJSS}, so non-degeneracy is the
familiar condition that the M\"uger center is trivial. Over an algebraically closed field
of characteristic zero, the three conditions of Definition \ref{defn:inv_E2} are mutually equivalent
for finite braided tensor categories, and a finite braided tensor category is invertible
if and only if it is non-degenerate \cite[Theorems 1.1 and 3.20]{BJSS}.

Examples include non-semisimple categories such as representation categories of small
quantum groups at suitable roots of unity, which give invertible fully extended framed
$4$-dimensional topological field theories in the abelian setting. Outside the finite setting this is not true:
The category of integrable representations of a quantum group at a generic
parameter is non-degenerate but not factorizable \cite[Remark 1.7]{BJSS}.

\subsection{Invertibility of commutative algebras}
\label{sec:inv_comm}

For a commutative algebra, invertibility in the higher Morita category is equivalent to
triviality.

\begin{prop}
\label{prop:inv_symm}
Let $A\in\CAlg\left(\coef\right)$ and $d\geq 1$.
Then $A$ is invertible in $\umor_d\left(\coef\right)$ if and only if the unit
$u\colon\one_{\coef}\to A$ is an equivalence.
\end{prop}

\begin{proof}
The ``if'' direction is clear, since then $A\simeq\one_{\coef}$ is the monoidal unit.
For ``only if'' we use the following observation.
Suppose we are given maps
$\one_{\coef}\xrightarrow{u}A\xrightarrow{i}W\xrightarrow{r}A$
with $ri\simeq\Id_A$ and $iu$ an equivalence.
Then $i$ is split epi, via $u\left(iu\right)^{-1}$, and split mono, via $r$, so $i$ is an
equivalence.
Hence $u\simeq i^{-1}\left(iu\right)$ is an equivalence.

Write $R_d\ceqq\int_{S^{d-1}\times\RR}A$, so that
$\mathcal{Z}_d\left(A\right)=\End_{R_d}\left(A\right)$.
If $A$ is invertible, then by Corollary \ref{cor:SSS_inv}, applied to $\int_{\emptyset}A=\one_{\coef}$, the unit
$\one_{\coef}\to\mathcal{Z}_d\left(A\right)$, which picks out $\Id_A$, is an equivalence.
Since $A$ is commutative, $R_d$ is the tensor $S^{d-1}\otimes A$ in $\CAlg\left(\coef\right)$
\cite[Section 5.5]{L:HA}, and it acts on $A=\int_{\RR^d}A$ through the map of commutative
algebras $\varepsilon\colon R_d\to A$ induced by $S^{d-1}\to\pt$.
Let $i\colon A\to\End_{R_d}\left(A\right)$ send $a$ to multiplication by $a$; this is $R_d$-linear
since $A$ is commutative and $R_d$ acts through $\varepsilon$.
By the unit axiom, $i\left(1_A\right)\simeq\Id_A$, so the unit of $\End_{R_d}\left(A\right)$ is $iu$.
Evaluation at $1_A$ defines $r\colon\End_{R_d}\left(A\right)\to A$ with $ri\simeq\Id_A$, and we
conclude with the above observation.
For $d=1$, $R_1=A\otimes A^{\rev}$ and this condition is centrality in \cite[Theorem 1.3]{BJSS}.
\end{proof}

Applying Proposition \ref{prop:inv_symm} with $\coef$ replaced by
$\Mod_R\left(\coef\right)$, whose unit is $R$, gives the relative version.

\begin{cor}
\label{cor:inv_symm_rel}
Let $R\in\CAlg\left(\coef\right)$, $A\in\CAlg\left(\Mod_R\left(\coef\right)\right)$ and
$d\geq 1$.
Then $A$ is invertible in $\umor_d\left(\Mod_R\left(\coef\right)\right)$ if and only if the
unit $R\to A$ is an equivalence.
\end{cor}

\subsection{Line operators in \texorpdfstring{$3$D}{three-dimensional} TQFTs}
\label{sec:line_ops}

Let $\cT$ be a symmetric monoidal $\left(\infty,3\right)$-category and
$\mathcal{F}\colon\Bord_3^{\mathrm{fr}}\to\cT$ a framed fully extended TQFT, with $B\ceqq
\mathcal{F}\left(\pt\right)$.
The \emph{category of line operators} of $\mathcal{F}$ is $\mathcal{F}\left(S^1\right)$, where $S^1$
carries the framing it inherits as the boundary of a disk; it is an $\EE_2$-algebra via embeddings
of disks.
By the cobordism hypothesis \cite{L:Cob},
\begin{equation*}
\mathcal{F}\left(S^1\right)\simeq\End_{\End_\cT\left(B\right)}\left(\Id_B\right) \ .
\end{equation*}
For $\cT=\umor_1\left(\Pre\right)$, $B$ is an $\EE_1$-algebra in $\Pre$,
$\End_\cT\left(B\right)$ is the category of $B$-bimodules and $\Id_B = B$. So the line
operators form the Drinfeld center
\begin{equation*}
\mathcal{F}\left(S^1\right)\simeq \mathcal{Z}_1\left(B\right)=\End_{B^e}\left(B\right)
\end{equation*}
of Notation \ref{nota:E2}.
Moreover $B$ is tautologically $\mathcal{Z}_1\left(B\right)$-central, i.e.\ it is a $1$-morphism
$\mathcal{Z}_1\left(B\right)\to\one$ in $\umor_2\left(\Pre\right)$, which is maximally adjointable since
$B$ is fully dualizable in $\umor_1\left(\Pre\right)$ (Remark \ref{rmk:SSS_relative}). So $\mathcal{F}$ is a
boundary condition for the $4$-dimensional theory determined by
$\mathcal{Z}_1\left(B\right)$ (when that theory exists).
